% Abridged version, Section 10: the finite computations that accompanied the work, and what remains conditional.
\section{Evidence and conditionality}
\label{sec:evi}
\noindent\genusstatus{numerically supported} The finite computations recorded in this section are exact or modular computations; they support the statements at the genera, orders and primes given, and are not proofs. The code and outputs are in \texttt{code/} of this project.


\emph{Finite computations.} Nothing in this paragraph and the next two is a step of a proof: the proofs of [L] are complete without it. We summarise the tables of \Lr{sec:evi}; a zero test modulo a prime is probabilistic. The outer terms $T_\Gamma=u^{2-2g}G_\Gamma(u)$, with $M=\Mo(\tau(u)\zh)$, were computed for every leg-free stable graph of genus $2$ to $5$, modulo $2^{89}-1$ and $2^{61}-1$ with rational reconstruction, in the frames $(\alpha,\beta)=(0,0)$ and $(1/3,-2/7)$. There are $7$, $42$, $379$ and $4555$ stable graphs; all $6$, $41$, $378$ and $4554$ edge graphs are regular (at genus $2$, $3$ and $5$ in the second frame modulo one prime), and the vertex graph, respectively the total, has the polar part $c_g\tc^{2-2g}$ with $c_g=\frac1{1200},\frac1{25200},\frac1{180000},\frac1{660000}$, in the frame $(0,0)$ and, for $g=3,4,5$, also in the second frame. At genus $4$ and $5$ the outer sum was obtained by regrouping the plain graph sum along the edges of $M$, a regrouping checked against genuine vertex tensors only at genus $2$ and $3$. As controls: with the bare symbol $\Mo(\zh)$ all $149$ edge graphs of genus $4$ with at most five edges are singular, and with $\tau$ replaced by $\tau+u^5$ exactly the three one-edge graphs of genus $4$ fail. The value $c_g$ and the gap of the total do not depend on $M$; only the regularity of the edge graphs tests the factorisation \Lr{sec:evi-graphs}.

The identity $\Dn_0=1$ was tested order by order at the conifold point up to $z^{204}$ (that is, $g\leq69$) exactly and up to $z^{1004}$ modulo two primes; the identity $(\Gamma)$, $c_n(1)=h_n$, for $n\leq60$ exactly and up to $z^{500}$ modulo two primes; $\PhiSP=\Phican$ up to $\zh^{80}$ and $u^{75}$ modulo primes; and $c(z;0)=\Ginf$ up to $z^{80}$ exactly \Lr{sec:evi-normal}.

For the model, \cref{thm:mod-tr}, the values of $\FTR_h$ for $h\leq5$, the lattice theorem and the Model Lemma were checked on finite ranges; off $\sK$ the law $\langle X\rangle_h=\Pi_0(X)(\Fm_h)'$ for singular one-point correlators, which holds on $\sK$ by $\ML{h}$, fails in an example \Lr{cor:mod-onepoint,sec:evi-model}.

\emph{What remains conditional.} Theorems~A and~C and Corollary~B are conditional on the inputs of \cref{sec:int-inputs}, and on nothing else: the quoted theorems (Q) on topological recursion, used for the model curve only, and on the action of $R$-matrices; the flatness of the $R$-matrix of the formal quintic and its value at the large-radius point (QRR); the genus-zero mirror theorem (M) and the genus-zero axioms (G0) of the twisted theory; Teleman's reconstruction with the axioms of the twisted theory in all genera (T); and the formula \cite[(1.7)]{AC96}. They are quoted with their statement numbers and are read with the conventions fixed in [L], where the sign conventions of the sources are compared \Lr{sec:sgn}. Theorem~Gamma and Theorem~U use no quoted input. Every other step is proved in [L]; in this abridged version the proofs are sketched, except where a proof is given in full from the stated lemmas. The limits of the results are those of \cref{sec:int-notclaimed}; in particular the conifold expansion is a formal change of frame, with no analytic continuation, and nothing is claimed for the quintic threefold itself.
